% Oscilador de 1 GDL: masa, muelle, amortiguador, fuerza y coordenada.
\documentclass[tikz,border=4pt]{standalone}
\usepackage{resonancia}
\begin{document}
\begin{tikzpicture}
  % apoyo
  \pared{(0,0)}{2.4}
  % masa sobre carro
  \node[masa, minimum width=2.0cm, minimum height=1.5cm, anchor=west] (M) at (4.0,0) {$m$};
  \draw[trazo] (4.3,-0.75) circle (0.14) (5.7,-0.75) circle (0.14);
  \suelo{(3.6,-0.89)}{(6.4,-0.89)}
  % muelle (arriba) y amortiguador (abajo)
  \resorte{(0,0.45)}{(4.0,0.45)}
  \amortiguador{(0,-0.35)}{(4.0,-0.35)}
  \node[etiqueta, above] at (2.0,0.78) {$k$};
  \node[etiqueta, below] at (2.0,-0.75) {$c$};
  % fuerza y coordenada
  \draw[fuerza] (6.0,0) -- (7.3,0) node[midway, above=1pt, etiqueta, color=rNaranja] {$f(t)$};
  \draw[eje] (4.0,1.45) -- (5.3,1.45) node[right, etiqueta] {$x(t)$};
  \draw[punteada] (4.0,0.85) -- (4.0,1.6);
  % ecuación
  \node[etiqueta, anchor=north] at (3.4,-1.5) {$m\ddot{x}+c\dot{x}+kx=f(t)$};
\end{tikzpicture}
\end{document}
